\section{Statement of the Problem}
Hospitals are critical facilities expected to remain available after earthquakes, yet structural and foundation deformations can interrupt medical services even when global collapse is avoided \cite{guragain2004,subedi2022}. In Nepal, high seismic hazard occurs together with variable soil conditions; consequently, a fixed-base idealization may not represent the response of a hospital founded on flexible soil with sufficient accuracy.

SSI changes the dynamic characteristics of the coupled soil--foundation--structure system and can redistribute seismic demand. These changes may appear in modal period, storey displacement, inter-storey drift, base and storey shear, foundation settlement, support reactions, and member forces. They may also influence nonlinear capacity and the intensity at which drift-based damage states are reached. Nepal-focused studies that examine these response measures together for a mid-size hospital building remain limited.

The problem addressed in this thesis is therefore the extent to which fixed-base and flexible-base idealizations lead to different seismic response and preliminary fragility estimates for the selected hospital prototype under hard, medium, and soft soil conditions.
